\vspace{-3.5mm}
\subsection{Overview of \sys} 
\vspace{-1mm}
\label{subsec:overview}

\rone{Our notion of provenance considers a ``black-box'' LLM that performs an unknown computation on a text data source whose output we observe. The goal of our provenance-finding algorithm is to identify the minimal text that can lead to an equivalent output. Because the underlying computation is unknown, our approach, unlike prior work on provenance in relational  databases~\cite{green2007provenance,buneman2001and,cheney2009provenance}, cannot track provenance throughout query evaluation. Instead, we develop a method that \textit{reverse-engineers} the LLM's answering process after the fact.   
}

We next present an overview of \sys, our
provenance discovery framework, as shown in  Figure~\ref{fig:overview}. 
Given text $T$, question $Q$, LLM $L$, and
answer $A = L(Q,T)$, 
\sys uses a two-phase approach to infer a single minimal provenance (Section~\ref{sec:1-provenance}),
and can be extended to find $k$ distinct minimal provenances  (Section~\ref{sec:k-provenance}). 


\topic{Finding a single minimal provenance}
\sys uses a two-phase approach, 
{\em prune-refine}, 
to return a minimal provenance. 
The {\em pruning} phase returns a small-sized provenance by aggressively pruning out sentences 
that likely do not contribute to $A$, 
while the {\em refinement} phase 
takes the provenance 
returned from the pruning phase 
and refines it to obtain a minimal provenance. 

For the pruning phase, 
we propose four strategies 
to reduce provenance size 
by leveraging {\em text relevance to $Q$ and $A$} and {\em different scan orders} (top-down or bottom-up).  
\vspace{-1mm}
\begin{itemize}[leftmargin=*]
    \item (Text Relevance) We divide the text $T$ into equal-sized portions and prioritize examining those with high relevance to the pair $\langle Q,A \rangle$, estimated using embedding similarity~\cite{caspari2024beyond} or an LLM. 
    \item (Scan Orders) To return a small-sized provenance quickly, one approach is to process text portions in decreasing order of relevance to $\langle Q,A \rangle$, incrementally add a portion at a time, and return the top‑$i$ portions as soon as they yield an answer equivalent to $A$.    
    Alternatively, starting from $T$, we repeatedly discard the least relevant half until the remaining portions fail to reproduce $A$. 
\end{itemize}
\vspace{-1mm}
As bottom-up and top-down approaches excel in different scenarios, we provide a theoretical cost analysis to characterize where each shines, based on the size of the returned provenance. We further propose an adaptive method that combines them to robustly infer a small-sized provenance while avoiding their worst-case limitations. 
%\agp{mention cost model here} 

In the refinement phase, we propose two strategies that take any provenance (e.g., the output from the prune phase) 
and  return a corresponding 
minimal provenance by removing 
all non-relevant sentences 
with respect to $\langle Q, A\rangle$. 
The first strategy, {\em sequential-greedy}, 
iteratively deletes sentences 
from a given provenance 
if the remaining text after each removal 
still produces an answer equivalent to $A$. 
In contrast to {\em sequential-greedy} 
that considers one sentence at a time, the second strategy, 
{\em exponential-greedy}, removes sentences in exponentially increasing batch sizes. Starting with one sentence, if deletion succeeds, it attempts to remove two, then four, and so on, doubling the batch size each time to accelerate pruning. Both strategies enumerate sentences in decreasing order of their indexes (from larger to smaller) to maximize token reuse via the KV cache in each LLM invocation, leading to a significant cost reduction. 
%\agp{constant reduction? i don't follow}

%Combining the above considerations, we conduct an exhaustive explorations by proposing four different strategies in the prune phase, and two strategies in the refine phase to return a minimal provenance.  Different combinations of strategy may perform better in different scenarios, depending on factors such as document and provenance size, question type. We present empirical observations and provide recommendations for choosing strategies under various settings in Section~\ref{sec:exp-top-1}. 


\topic{Find top-$k$ minimal provenances}
The provenance returned by the two-phase approach
is minimal. However, 
in certain scenarios, as outlined in Section~\ref{sec:k-provenance}, additional evidence is needed. 
Thus, as part of \sys, we introduce a 
greedy tree-search approach 
to return the $k$ distinct 
minimal provenances, 
allowing users to explore all provenances and gain greater confidence in the answers. 
%\agp{notion of validity has not been introduced. i would 
%play it safe and just say we return $k$ distinct provenances,
%not top-k.
%}

%However, it may not be valid. (e.g., the third provenance in Figure~\ref{fig:example}-1 is minimal but invalid.) 







 